\ExerciseSection

\begin{enumerate}
\def\labelenumi{\arabic{enumi})}
\tightlist
\item
  \begin{enumerate}
  \def\labelenumii{\alph{enumii})}
  \tightlist
  \item
    设 A 是拓扑空间 X 的子集，试证下列命题等价：\(\alpha\) A 的边界无处稠密；\(\beta\) A 是一个开集与一个无处稠密集的并；\(\gamma\) A 是一个闭集与一个无处稠密集的差。
  \end{enumerate}
\end{enumerate}

\begin{enumerate}
\def\labelenumi{\alph{enumi})}
\setcounter{enumi}{1}
\tightlist
\item
  试证：若 X 的子集 A 使得对每一点 \(x \in A\)，存在 x 在 X 中的邻域 V 满足 \(V \cap A\) 在 X 中无处稠密，则 A 在 X 中无处稠密。
\end{enumerate}

\begin{enumerate}
\def\labelenumi{\arabic{enumi})}
\setcounter{enumi}{1}
\tightlist
\item
  设 \(\mathbf{A}\) 是拓扑空间 \(\mathbf{X}\) 的子集，称它为瘦集，如果对每个完备集 \(\mathbf{P} \subset \mathbf{X}\)，\(\mathbf{P} \cap \mathbf{A}\) 在 \(\mathbf{P}\) 中无处稠密。
\end{enumerate}

\begin{enumerate}
\def\labelenumi{\alph{enumi})}
\item
  瘦集的有限并是瘦集。
\item
  瘦集的边界无处稠密。
\item
  \(\mathbf{X}\) 中存在一个最大的完备集，其补集是瘦集。
\end{enumerate}

\begin{enumerate}
\def\labelenumi{\arabic{enumi})}
\setcounter{enumi}{2}
\tightlist
\item
  设 A 是拓扑空间 X 的子集，用 D(A) 表示满足下列条件的所有 \(x \in X\) 组成的集：对 x 的每个邻域 V，集合 \(V \cap A\) 不是贫集。D(A) 包含于 \(\overline{A}\)。
\end{enumerate}

\begin{enumerate}
\def\labelenumi{\alph{enumi})}
\tightlist
\item
  若 \(\mathbf{B}\) 是 \(\mathbf{X}\) 的子集且 \(\mathbf{D}(\mathbf{B}) = \varnothing\)，则
\end{enumerate}

\[
\mathbf {D} (\mathbf {A} \cup \mathbf {B}) = \mathbf {D} (\mathbf {A})
\]

对 X 的每个子集 A 成立。

\begin{enumerate}
\def\labelenumi{\alph{enumi})}
\setcounter{enumi}{1}
\item
  D(A) = ∅ 当且仅当 A 是贫集。{[}先证明：若 (U\_t) 是 X 的两两不相交的开子集族，且对每个 t，V\_t 是包含于 U\_t 中的（相对于 X 的）无处稠密集，则 ∪ V\_t 无处稠密。然后考虑 X 的开子集的一个极大集 M，其中每两个开集都不相交，且对每个 U ∈ M，A ∪ U 都是贫集；这样的极大集的存在性由佐恩引理建立。证明若 D(A) = ∅，则 A ∩ C W 无处稠密，其中 W 是 M 中诸集的并，进而证明 A 是贫集。{]}
\item
  试证 \(\mathbf{D}(\mathbf{A})\) 是闭集，且 \(\mathbf{A} \cap \mathbf{CD}(\mathbf{A})\) 对所有 \(A \subset X\) 是贫集{[}借助 b)，证明 \(\mathbf{A} \cap \mathbf{CD}(\mathbf{A})\) 相对于开子空间 \(\mathbf{CD}(\mathbf{A})\) 是贫集{]}。
\end{enumerate}

\(d\) ) 试证 \(\mathbf{D}(\mathbf{A})\) 等于其内部的闭包。{[}若 \(\mathbf{D}'(\mathbf{A})\) 是 \(\mathbf{D}(\mathbf{A})\) 的内部的闭包，注意 \(\mathbf{A} \cap \mathbb{C}\mathbf{D}'(\mathbf{A})\) 是以下两集之并：

\[
\mathbf {A} \cap \mathbb {C D} (\mathbf {A}) \quad \text { and } \quad \mathbf {A} \cap \mathbf {D} (\mathbf {A}) \cap \mathbb {C D} ^ {\prime} (\mathbf {A}). ]
\]

\begin{enumerate}
\def\labelenumi{\arabic{enumi})}
\setcounter{enumi}{3}
\tightlist
\item
  设 \(\mathbf{X}\)、\(\mathbf{Y}\) 是两个拓扑空间，\(\mathbf{A}\)（相应地，B）是 \(\mathbf{X}\)（相应地，Y）的子集。
\end{enumerate}

\begin{enumerate}
\def\labelenumi{\alph{enumi})}
\item
  \(\mathbf{A} \times \mathbf{B}\) 在 \(\mathbf{X} \times \mathbf{Y}\) 中无处稠密（相应地，完备）当且仅当两集 \(\mathbf{A}, \mathbf{B}\) 之一无处稠密（相应地，\(\mathbf{A}\) 与 \(\mathbf{B}\) 均为闭集且其中之一完备）。
\item
  \(\mathbf{A} \times \mathbf{B}\) 在 \(\mathbf{X} \times \mathbf{Y}\) 中是瘦集当且仅当 \(\mathbf{A}\) 与 \(\mathbf{B}\) 都是瘦集。
\end{enumerate}

\begin{enumerate}
\def\labelenumi{\arabic{enumi})}
\setcounter{enumi}{4}
\tightlist
\item
  \begin{enumerate}
  \def\labelenumii{\alph{enumii})}
  \tightlist
  \item
    设 X、Y 是两个拓扑空间，A 是 \(X \times Y\) 的无处稠密子集。若 Y 的拓扑具有可数基，试证由所有 \(x \in X\) 中使 A 在 x 处的截面 \(\mathbf{A} \cap (\{x\} \times \mathbf{Y})\) 相对于 \(\{x\} \times Y\) 不是无处稠密者组成的集是 X 中的贫集（若 U 是 Y 中任一非空开集，试证由所有 \(x \in X\) 中使 \(\{x\} \times U\) 包含于 A 在 x 处截面的闭包之中者组成的集是 X 中的无处稠密集）。
  \end{enumerate}
\end{enumerate}

\begin{enumerate}
\def\labelenumi{\alph{enumi})}
\setcounter{enumi}{1}
\item
  举出一个例子说明，若 Y 的拓扑不具有可数基，则 a) 的结论可以不成立（取 X 为无孤立点的豪斯多夫空间，Y 为赋予离散拓扑的集合 X；并取 A 为 \(X \times Y\) 的对角线）。
\item
  若 B、C 分别是 X、Y 的子集，且 B、C 之一（分别在 X、Y 中）是贫集，则 \(B \times C\) 在 \(X \times Y\) 中是贫集。反之，若 \(B \times C\) 在 \(X \times Y\) 中是贫集，且 X 或 Y 之一的拓扑具有可数基，则 B、C 两集之一是贫集。
\item
  由 c) 推出：若空间 X、Y 之一的拓扑具有可数基，则在习题 3 的记号下有 \(\mathrm{D}(\mathbf{B} \times \mathbf{C}) = \mathrm{D}(\mathbf{B}) \times \mathrm{D}(\mathbf{C})\)。
\end{enumerate}

\begin{enumerate}
\def\labelenumi{\arabic{enumi})}
\setcounter{enumi}{5}
\tightlist
\item
  拓扑空间 X 中的集合 A 称为几乎开集，如果存在 X 中的开集 U，使 \(U \cap C A\) 与 \(A \cap C U\) 都是贫集。
\end{enumerate}

\begin{enumerate}
\def\labelenumi{\alph{enumi})}
\item
  几乎开集的补集是几乎开集。
\item
  几乎开集的任意可数并是几乎开集。
\item
  试证下列命题等价：\(\alpha\) ) A 是几乎开集；\(\beta\) ) 存在 X 的贫子集 M，使 A∩CM 在子空间 X---M 中既开又闭；\(\gamma\) ) 存在集合 G⊂A，它是 X 中开集的可数交，且 A---G 在 X 中是贫集；\(\delta\) ) 存在集合 F⊃A，它是 X 中闭集的可数并，且 F---A 在 X 中是贫集；\(\zeta\) ) 集合 D(A)∩D(X---A)（习题 3）在 X 中无处稠密；\(\theta\) ) 集合 D(A)∩CA 在 X 中是贫集。{[}利用习题 3a) 与 3b) 证明 \(\alpha)\Longrightarrow\zeta)\Longrightarrow\theta)\Longrightarrow\alpha\)。{]} d) 设 X、Y 是两个拓扑空间，且 X 或 Y 之一的拓扑具有可数基。则 X × Y 中形如 \(\mathbf{A} \times \mathbf{B}\) 的子集是几乎开集，当且仅当 \(\mathbf{A}\) 与 \(\mathbf{B}\) 都是几乎开集，或者其中之一是贫集{[}利用 \(\zeta\) )，即 \(c\) 中的相应断言{]}。
\end{enumerate}

\begin{enumerate}
\def\labelenumi{\arabic{enumi})}
\setcounter{enumi}{6}
\tightlist
\item
  拓扑空间 X 称为非贫的，如果它相对于自身不是贫集。
\end{enumerate}

\begin{enumerate}
\def\labelenumi{\alph{enumi})}
\item
  X 非贫当且仅当 X 中稠密开集的每一可数族具有非空交。
\item
  在非贫空间中，每个贫集的补集都是非贫子空间。
\end{enumerate}

\begin{enumerate}
\def\labelenumi{\arabic{enumi})}
\setcounter{enumi}{7}
\item
  设 X 是一个拓扑空间，其中存在非空开子集 A，它作为子空间是非贫的。试证 X 非贫，且 A 相对于 X 不是贫集。
\item
  设 X 是 \(R^{2}\) 的子空间，它是 \(Q^{2}\) 与直线 y = 0 的并。试证 X 相对于自身是贫集。由此推出：一个拓扑空间可以有非贫的子空间，而其自身却不是非贫的。
\item
  试证每个完全正则的伪紧空间（§ 1，习题 21）都是贝尔空间（按定理 1 的方式论证）。
\item
  若 X 是无孤立点的贝尔空间，试证 X 的每一点都是 X 的凝聚点（第 I 章，§ 9，习题 17）。
\end{enumerate}

¶ 12) 拓扑空间 X 称为完全非贫的，如果 X 的每个非空闭子空间都非贫。局部紧空间是完全非贫的，完备度量空间也是如此。

\begin{enumerate}
\def\labelenumi{\alph{enumi})}
\item
  试证完全非贫的正则空间是贝尔空间。
\item
  在可达的完全非贫空间中，非空可数闭集是瘦集（习题 2），因而至少有一个孤立点。
\item
  在完全非贫空间 X 中，每个作为 X 中开集的可数交的子空间 A 都是完全非贫的（设 F 是 A 的在 A 中闭的子集，用 \(\overline{F}\) 表示 F 在 X 中的闭包，并证明 \(\overline{F} \cap C F\) 相对于 \(\overline{F}\) 是贫集）。
\item
  若拓扑空间 X 的每一点都有一个作为完全非贫子空间的邻域，则 X 是完全非贫的。
\end{enumerate}

¶ 13) 设 X 是连通且局部连通的完全非贫空间。试证 X 不能是非空互不相交闭集的无穷序列 (F \(_{n}\) ) 的并。（设 H 是诸 F \(_{n}\) 在 X 中的边界 H \(_{n}\) 的并；证明 H 在 X 中闭，且每个 H \(_{n}\) 在 H 中无处稠密；为证明后一点，考虑 \(\mathbf{H}_n\) 的点的基本邻域系（由连通邻域组成），并用反证法，利用第 I 章 § 11 命题 3。）

\begin{enumerate}
\def\labelenumi{\arabic{enumi})}
\setcounter{enumi}{13}
\tightlist
\item
  设 X 是 \(R^{2}\) 的子空间，由点 \((r,0)\)（其中 r 取遍有理数集 Q）与点 \((k/n,1/n)\)（其中 n 取遍整数集 \(\geqslant 1\)，k 取遍全体有理整数）组成。试证 X 是贝尔空间但不是完全非贫的。
\end{enumerate}

¶ 15) 设 X 是平面 \(\mathbf{R}^2\) 的子集，由直线 \(\mathrm{D} = \{0\} \times \mathbf{R}\) 与点 \((1/n, k/n^2)\)（其中 \(n\) 取遍整数集 \(>0\)，\(k\) 取遍全体有理整数）组成。a) 对 D 的每一点 \((0, y)\) 与每个整数 \(n > 0\)，用 \(\mathrm{T}_n(y)\) 表示由点 \((u, v) \in \mathbf{X}\)（满足 \(u \leqslant 1/n\) 且 \(|v - y| \leqslant u\)）组成的集。试证：若取每一点 \((0, y)\) 的由诸集 \(\mathrm{T}_n(y)\) 组成的集作为基本邻域系，并取由该点单独组成的集作为 X 的其余每一点的基本邻域系，则我们在 X 上定义了一个拓扑 \(\mathcal{G}\)，它比 \(\mathbf{R}^2\) 的拓扑在 X 上诱导的拓扑更细，并且 X 关于它是局部紧的。b) 试证 X 是可数个可度量化闭子空间的并，且 X 的每一点都有一个可度量化的紧邻域，但 X 不是正规空间。{[}设 A 是由所有点 \((0, y)\)（其中 \(y\) 为有理数）组成的集，并设 \(\mathbf{B} = \mathbf{D} - \mathbf{A}\)；证明 A 在 X 中的每个邻域 U 都与 B 在 X 中的每个邻域 V 相交；为此，对每个 \(n\)，考察集 \(\mathbf{B}_n\)，即由这样的点 \(y \in \mathbf{B}\) 组成的集：\(\mathrm{T}_n(y) \subset \mathbf{V}\)；并注意集 B 在 R 中不是贫集。{]}

¶ 16) 用 \(\mathbf{R}_{-}\) 表示在实数的线性序集 \(\mathbf{R}\) 上赋予拓扑 \(\mathcal{G}_{-}(\mathbf{R})\)（第 I 章，§ 2，习题 5）所得的拓扑空间。

\begin{enumerate}
\def\labelenumi{\alph{enumi})}
\item
  试证 \(\mathbf{R}_{-}\) 中的每个开集都是可数个互不相交区间之并，这些区间或者是开区间，或者是左半开区间。由此推出 \(\mathbf{R}_{-}\) 是完备正规的（§ 4，习题 8）且是完全非贫的（习题 12）。
\item
  试证空间 \(\mathbf{R}_{-} \times \mathbf{R}_{-}\) 不是正规空间，从而 \(\mathbf{R}_{-}\) 不是可度量化的，尽管 \(\mathbf{R}_{-}\) 含有可数稠密子集。（考虑由点 \((x, y)\)（在 \(\mathbf{R}_{-} \times \mathbf{R}_{-}\) 中且满足 \(x + y = 1\)）组成的集，证明 \(\mathbf{R}_{-} \times \mathbf{R}_{-}\) 有一个闭子空间同胚于习题 15 中定义的空间 \(\mathbf{X}\)。）
\item
  试证 \(\mathbf{R}_{-}\) 是林德勒夫空间（第 I 章，§ 9，习题 15），因而是仿紧的（§ 4，习题 23）。（对每个 \(x \in \mathbf{R}_{-}\)，设 \(\mathbf{U}_x\) 是半开区间 {]} \(y(x)\)，\(x]\)，以 \(x\) 为右端点。试证由所有 \(z \in \mathbf{R}_{-}\) 中不属于任何开区间 {]} \(y(x)\)，\(x[\) 者组成的集是可数的；利用第 IV 章 § 2 习题 1。）
\item
  试证 \(R_{-}\) 的每个紧子集都是可数的{[}若 A 是 \(R_{-}\) 中的紧集，则 \(\mathcal{G}_{-}(\mathbf{R})\) 与 \(\mathcal{G}_{0}(\mathbf{R})\) 在 A 上诱导的拓扑相同；由此推出 A 的每一点都是 A 的邻接区间的左端点{]}。
\end{enumerate}

\begin{enumerate}
\def\labelenumi{\arabic{enumi})}
\setcounter{enumi}{16}
\tightlist
\item
  \begin{enumerate}
  \def\labelenumii{\alph{enumii})}
  \tightlist
  \item
    试证完备度量空间的每个乘积 \(\mathbf{X} = \prod_{i\in \mathbf{I}}\mathbf{X}_i\) 都是贝尔空间。{[}按定理 1 的证法论证，取 \(\mathbf{G}_n\) 为初等集（第 I 章，§ 4，第 1 节），其所有投影或为整个空间，或直径 \(\leqslant 1 / n\)；然后注意存在 I 的可数子集 J，使每个 \(\mathbf{G}_n\) 形如 \(\mathbf{H}_n\times \prod_{i\notin \mathbf{J}}\mathbf{X}_i\)，其中 \(\mathbf{H}_n\subset \prod_{i\in \mathbf{J}}\mathbf{X}_i.\){]}
  \end{enumerate}
\end{enumerate}

\begin{enumerate}
\def\labelenumi{\alph{enumi})}
\setcounter{enumi}{1}
\tightlist
\item
  举出一个例子说明：若 I 不可数，则 X 未必是完全非贫的（习题 12）。{[}利用 § 2 的习题 1 与 23 b)。{]}
\end{enumerate}

¶ 18) a) 试证在完备度量空间中，每个完备集都含有一个同胚于康托尔三分集（第 IV 章，§ 2，第 5 节）的子集（利用第 IV 章 § 8 的习题 11）。

\begin{enumerate}
\def\labelenumi{\alph{enumi})}
\setcounter{enumi}{1}
\item
  由此推出：在具有可数基的完备度量空间中，每个完备集都具有连续统的势，且每个闭集与每个开集或者是可数的，或者具有连续统的势{[}利用习题 11 b){]}。
\item
  试证在具有可数基的不可数完备度量空间中，完备子集的集具有连续统的势。{[}先对紧空间 \(\{0, 1\}^{N}\) 证明该断言，再利用 a) 与 § 2 习题 12 b)。{]}
\item
  设 X 是具有可数基的不可数完备度量空间。试证存在 X 的子集 Z，使 Z 与 X---Z 都具有连续统的势，且 Z 与 X---Z 都不含非空完备集。{[}利用 b)、c) 以及《集合论》第 III 章 § 6 习题 24 a) 中描述的方法。{]}进而，若 X 无孤立点，试证 Z 与 X---Z 都不是贫集。（若 A 是 X 的稠密开子集的可数交，试证 A 含有非空完备集，利用 c) 与 § 2 习题 24。）由此推出，在此情形 Z 不是 X 中的几乎开集（习题 6；注意上面的论证表明，若 U 是 X 的任一非空开子集，则 U ∩ Z 不是贫集）。
\end{enumerate}

\begin{enumerate}
\def\labelenumi{\arabic{enumi})}
\setcounter{enumi}{18}
\tightlist
\item
  设 A 是拓扑空间 X 的稠密子空间，f 是 A 到完备度量空间 Y 的连续映射。试证 X 中使 f（相对于 A）没有极限的点组成的集是 X 中的贫集{[}对每个 n \textgreater{} 0，考察点 \(x \in E\) 中使 f 的振幅 \(\omega(x; f)\) \textless{} 1/n 者组成的集{]}。
\end{enumerate}

¶ 20) 设 \(f\) 是完备度量空间 X 的子集 A 到子集 \(\mathbf{A}'\)（属于某个完备度量空间 \(\mathbf{X}'\)）上的同胚。试证存在 \(\overline{f}\)（\(f\) 到集 B 上的一个扩张，其中 B 是 X 中开集的可数交），使 \(f\) 是 B 到某个集 \(\mathbf{B}'\) 上的同胚，而该集是 \(\mathbf{X}'\) 中开集的可数交（对 \(f\) 与逆同胚 \(g\) 应用习题 19）。

¶ 21) a) 设 \((f_n)\) 是拓扑空间 X 到完备正规空间 Y（§ 4，习题 8）的连续映射序列，使得对每个 \(x \in \mathbf{X}\)，序列 \((f_n(x))\) 在 Y 中有极限 \(f(x)\)。试证对 Y 的每个闭子集 S，\(A = \overline{f}(S)\) 是开集的可数交。{[}存在 Y 中开集的递减序列 \((U_n)\)，使 \(S = \bigcap_{n} \overline{U}_n\)；注意 \(f(x) \in S\) 当且仅当对每个整数 \(n > 0\)，存在整数 \(k \geqslant 0\) 使 \(f_{n+k}(x) \in U_n\)。{]}由此推出 \(\overline{A} - A\) 在 X 中是贫集。

\begin{enumerate}
\def\labelenumi{\alph{enumi})}
\setcounter{enumi}{1}
\tightlist
\item
  由 a) 推出：若再假设 Y 是可数型的可度量化空间，则使 f 在其处不连续的所有 \(x \in X\) 组成的集是贫集。{[}若 \((\mathbf{V}_{n})\) 是 Y 的拓扑的基，且 \(S_{n} = Y - V_{n}\)，\(\mathbf{A}_{n} = f^{-1}(\mathbf{S}_{n})\)，试证 x 必属于诸集 \(\overline{\mathbf{A}}_{n} - \mathbf{A}_{n}\) 之一。{]}
\end{enumerate}

\begin{enumerate}
\def\labelenumi{\arabic{enumi})}
\setcounter{enumi}{21}
\tightlist
\item
  \begin{enumerate}
  \def\labelenumii{\alph{enumii})}
  \tightlist
  \item
    设 X 是贝尔空间，Y 是度量空间，\((f_{n})\) 是 X 到 Y 的连续映射序列，使得对每个 \(x \in X\)，序列 \((f_{n}(x))\) 在 Y 中有极限 \(f(x)\)。则使 f 不连续的所有点 \(x \in X\) 组成的集是贫集。{[}证明：对每个整数 n \textgreater{} 0，由点 \(x \in X\) 中使 f 的振幅 \(\omega(x; f)\) \(\leqslant 1/n\) 者组成的集包含一个稠密开集；为此，若用 \(G_{p}\) 表示由这样的 \(x \in X\) 组成的集：\(f_{p}(x)\) 到 \(f_{q}(x)\) 的距离 \(\leqslant 1/2n\)（对所有 \(q \geqslant p\)），试证诸集 \(G_{p}\) 的内部的并是一个稠密开集。{]}
  \end{enumerate}
\end{enumerate}

\begin{enumerate}
\def\labelenumi{\alph{enumi})}
\setcounter{enumi}{1}
\tightlist
\item
  设 K 是 R 中的紧区间 {[}0, 1{]}。给出连续映射序列 \(f_{n}\)（从 K 到紧空间 \(K^{K}\)）的一个例子，使得 \(\lim_{n\to\infty}f_{n}(x)=f(x)\) 对每个 \(x\in K\) 存在，但 f 在 K 的任何一点处都不连续。
\end{enumerate}

¶ 23) 设 X 是拓扑空间，Y、Z 是两个度量空间，其度量分别为 d、d'。设 f: X × Y → Z 是这样一个映射：对每个 x₀ ∈ X，y → f(x₀, y) 在 Y 上连续；对每个 y₀ ∈ Y，x → f(x, y₀) 在 X 上连续。

\begin{enumerate}
\def\labelenumi{\alph{enumi})}
\item
  对每个实数 \(\varepsilon > 0\)、每个 \(b \in \mathbf{Y}\) 与每个 \(x \in \mathbf{X}\)，用 \(g(x; b, \varepsilon)\) 表示满足下列条件的数 \(\alpha > 0\) 的上确界：由关系 \(d(b, y) < \alpha\) 可推出 \(d'(f(x, b), f(x, y)) \leqslant \varepsilon\)。试证函数 \(x \to g(x; b, \varepsilon)\) 是上半连续的。
\item
  若 X 是贝尔空间，试由 a) 与定理 2 推出：对每个 \(b \in Y\)，存在 X 中的集 \(S_{b}\)，其补集是贫集，使得 f 在 \((a, b)\) 处对每个 \(a \in S_{b}\) 连续。
\end{enumerate}

\begin{enumerate}
\def\labelenumi{\arabic{enumi})}
\setcounter{enumi}{23}
\tightlist
\item
  设 X、Y 是两个拓扑空间。映射 \(f: \mathbf{X} \to \mathbf{Y}\) 称为几乎开的，如果对 Y 的每个闭子集 S，\(\overline{f}(\mathbf{S})\) 在 X 中是几乎开集（习题 6）。
\end{enumerate}

\begin{enumerate}
\def\labelenumi{\alph{enumi})}
\item
  设 \(g: \mathbf{X} \to \mathbf{Y}\) 是具有下列性质的映射：存在贫集 \(\mathbf{M}\)（含于 \(\mathbf{X}\)），使 \(g|_{\mathbf{CM}}\) 连续。试证 \(g\) 是几乎开的。在 \(\mathbf{Y}\) 的拓扑具有可数基的情形，考察其逆命题。
\item
  假设 Y 是完备正规的（§ 4，习题 8）。设 \((f_n)\) 是 X 到 Y 的几乎开映射序列，使得
\end{enumerate}

\[
\lim _ {n \rightarrow \infty} f _ {n} (x) = f (x)
\]

对每个 \(x \in \mathbf{X}\) 存在。试证 \(f\) 是几乎开的{[}按习题 21 a) 的方式论证{]}。

\begin{enumerate}
\def\labelenumi{\arabic{enumi})}
\setcounter{enumi}{24}
\tightlist
\item
  设 R 是拓扑空间 X 上的开等价关系。试证：若 X 是非贫的（习题 7）{[}相应地，是贝尔空间；相应地，是完全非贫的（习题 12）{]}，则 X/R 亦然。
\end{enumerate}

¶ 26) 设 X 是局部紧的可度量化空间，d 是与 X 的拓扑相容的度量，R 是 X 上的闭等价关系，其所有等价类都是紧的。则 X/R 是局部紧的（第 I 章，§ 10，命题 9）且可度量化（§ 4，习题 24）。

\begin{enumerate}
\def\labelenumi{\alph{enumi})}
\tightlist
\item
  对每一对点 \(u, v\)（属于 \(\mathbf{X} / \mathbf{R}\)），令
\end{enumerate}

\[
\rho (u, v) = \sup _ {\varphi (x) = u} d (x, \overline {{\varphi}} ^ {1} (v)),
\]

其中 \(\varphi : \mathbf{X} \to \mathbf{X} / \mathbf{R}\) 是典范映射，且对每个 \(u \in \mathbf{X} / \mathbf{R}\)，令

\[
\lambda (u) = \lim _ {v \rightarrow u, v \neq u} \sup _ {\rho (u, v)}.
\]

试证 \(\lambda\) 在 X/R 上是上半连续的，且对每个 \(\alpha > 0\)，由点 \(u \in \mathrm{X} / \mathrm{R}\) 中使 \(\lambda(u) \geqslant \alpha\) 者组成的集是无处稠密的。{[}否则将存在 X 的紧子集 K 与 K 中点的无穷序列 \((x_n)\)，使 \(d(x_m, x_n) > \frac{1}{2} \alpha\)（对 \(m \neq n\)）。{]} b) 由 a) 推出：存在 X/R 的稠密子集 M，它是 X/R 中开集的可数交，且对每个 \(x \in \varphi^{-1}(M)\)，x 在 X 中的每个邻域在 \(\varphi\) 下的像是 \(\varphi(x)\) 在 X/R 中的邻域（利用定理 1）。

\begin{enumerate}
\def\labelenumi{\arabic{enumi})}
\setcounter{enumi}{26}
\tightlist
\item
  \begin{enumerate}
  \def\labelenumii{\alph{enumii})}
  \tightlist
  \item
    设 G 是拓扑群，A 是 G 中的几乎开集（习题 6）。试证：若 A 不是贫集，则 AA \(^{-1}\) 是 G 中单位元的邻域。{[}首先注意，由习题 3 知 G 是贝尔空间。设 A* 是所有满足 U ∩ C A 为贫集的开集 U ⊂ G 的并。证明 A* 非空，且若 x ∈ G 使 xA* 与 A* 相交，则 xA 与 A 相交。由此证明 AA \(^{-1}\) ⊃ A*(A*) \(^{-1}\)。{]}
  \end{enumerate}
\end{enumerate}

\begin{enumerate}
\def\labelenumi{\alph{enumi})}
\setcounter{enumi}{1}
\item
  由 a) 推出：拓扑群 G 的几乎开子群或者是贫集，或者是开的（且闭）。特别地，作为豪斯多夫贝尔空间的可数拓扑群是离散的。
\item
  试证拓扑群 R 存在这样的子群 H，使 R/H 是可数无穷的（利用哈梅尔基）；这样的子群是稠密的，但既不是贫集也不是几乎开集。
\item
  设 G 是左、右一致结构一致的可度量化拓扑群。试证：若在 G 上存在与 G 的拓扑相容的度量，使 G 关于它是完备度量空间，则 G 是完备可度量化群{[}利用 b)；§ 3 习题 2；以及 § 2 习题 24{]}。
\end{enumerate}

\begin{enumerate}
\def\labelenumi{\arabic{enumi})}
\setcounter{enumi}{27}
\tightlist
\item
  设 \(\mathbf{G}, \mathbf{G}'\) 是两个拓扑群。试证每个连续同态 \(f\)（\(\mathbf{G}\) 到 \(\mathbf{G}'\) 上的）在下列两种情形的每一种中都是严格态射：
\end{enumerate}

\begin{enumerate}
\def\labelenumi{\alph{enumi})}
\item
  G 局部紧且 \(\sigma\) 紧，而 \(\mathbf{G}'\) 非贫{[}证明 G 中存在相对紧的开集 U，使 \(f(\overline{\mathbf{U}})\) 的内部非空；由此推出，对 G 的单位元的每个紧邻域 V，\(f(\mathbf{V})\) 的内部非空，然后应用第 III 章 § 2 习题 18{]}。
\item
  G 完备，G 的拓扑具有可数基，且 \(G'\) 非贫。{[}利用如下事实：对 G 的单位元 e 的每个邻域 U，G 都是形如 \(x_{n}U\) 的集的序列之并，证明对 G 的每个开集 A，存在开集 \(A'\)（在 \(G'\) 中），它包含 \(f(A)\) 且使 \(f(A)\) 在 \(A'\) 中稠密。然后考虑 e 在 G 中的对称邻域的基本系 \((\mathbf{U}_{n})\)，使 \(\overline{U}_{n+1}^{2} \subset U_{n}\)，并证明 \(f(\mathbf{U}_{p})\) 包含开集 \(U_{p+1}'\)；为此，取任一点 \(a' \in U_{p+1}'\)，然后用归纳法构造 G 中点的序列 \((b_{n})\)，使 \(b_{n} \in U_{p+n}\) 且 \(f(b_{1}b_{2}\ldots b_{n})\) 趋于 \(a'\)；注意，若 \(x' \in U_{k}'\)，则存在 \(y \in \mathbf{U}_k\) 使 \(f(y) \in x' \mathrm{U}_{k+1}'\)。{]}关于 \(\mathbf{G}\) 具有可数基的假设是本质的，如下例所示：\(\mathbf{G}'\) 是带实直线拓扑的群 \(\mathbf{R}\)，\(\mathbf{G}\) 是带离散拓扑的群 \(\mathbf{R}\)，而 \(f: \mathbf{G} \to \mathbf{G}'\) 是恒同映射。
\end{enumerate}

\begin{enumerate}
\def\labelenumi{\arabic{enumi})}
\setcounter{enumi}{28}
\tightlist
\item
  设 G 是局部紧、σ 紧的拓扑群，或者是拓扑具有可数基的完备可度量化群。设 H 是 G 的闭正规子群，A 是 G 的闭子群。则商群 A/(A∩H) 同构于 AH/H 当且仅当 AH 是 G 的闭子群{[}为证条件必要，注意由第 III 章 § 4 第 6 号命题 13 与第 IX 章 § 3 第 1 号命题 4 知 A/(A∩H) 完备；为证条件充分，利用习题 28{]}。
\end{enumerate}

¶ 30) 设 G 是群，d 是 G 上的度量，使得关于 d 在 G 上定义的拓扑，G 到自身的映射 \(y \to x_{0}y\) 与 \(y \to yx_{0}\) 对每个 \(x_{0} \in G\) 都连续。

\begin{enumerate}
\def\labelenumi{\alph{enumi})}
\item
  试证：若 \(\mathbf{G}\) 关于拓扑 \(\mathfrak{C}\) 完备，则映射 \((x,y)\to xy\) 在 \(\mathbf{G}\times\mathbf{G}\) 上连续（利用习题 23）。
\item
  再假设 \(\mathcal{C}\) 具有可数基。试证映射 \(x \to x^{-1}\) 在 \(\mathbf{G}\) 上连续，即拓扑 \(\mathcal{C}\) 与 \(\mathbf{G}\) 的群结构相容。{[}在群 \(\mathbf{G} \times \mathbf{G}\)（赋予拓扑 \(\mathcal{C} \times \mathcal{C}\)）中，考察集 \(\mathbf{F}\)，它由所有点 \((x, x^{-1})\)（\(x \in \mathbf{G}\)）组成；\(\mathbf{F}\) 在 \(\mathbf{G} \times \mathbf{G}\) 中闭，且合成律 \((x, x^{-1})(y, y^{-1}) = (xy, y^{-1}x^{-1})\) 在 \(\mathbf{F}\) 上定义了一个群结构；试证 \(\mathbf{F}\) 上由 \(\mathbf{G} \times \mathbf{G}\) 的拓扑诱导的拓扑与此群结构相容（利用 \(a\)）。然后按习题 28 \(b\) 的论证方式，证明投影 \(\mathrm{pr}_1\)（\(\mathbf{F}\) 到 \(\mathbf{G}\) 上的）是双连续的。{]}
\end{enumerate}

